\DocumentMetadata{
  pdfversion=2.0,
  pdfstandard=ua-2,
  lang=en-US,
  tagging=on,
  tagging-setup={math/setup=mathml-SE,tabsorder=structure}
}
\documentclass[11pt,letterpaper]{article}
\usepackage[margin=0.68in,includefoot]{geometry}
\usepackage{newcomputermodern}
\usepackage{amsmath,amscd,mathtools}
\providecommand{\square}{\mdlgwhtsquare}
\usepackage[hidelinks]{hyperref}
\hypersetup{
  pdftitle={Topology First-Year Exam, May 2025, Part 1},
  pdfauthor={University of Florida Department of Mathematics},
  pdfsubject={Graduate mathematics examination},
  pdfdisplaydoctitle=true,
  bookmarksopen=true
}
\setcounter{secnumdepth}{0}
\setlength{\parindent}{0pt}
\setlength{\parskip}{0.28em}
\setlength{\emergencystretch}{3em}
\setlength{\leftmargini}{2.15em}
\setlength{\leftmarginii}{2.65em}
\setlength{\leftmarginiii}{3em}
\renewcommand{\labelenumi}{\textbf{\theenumi.}}
\pagestyle{empty}
\raggedbottom
\allowdisplaybreaks
\newenvironment{examproblems}[1]{%
  \begin{enumerate}%
  \setcounter{enumi}{#1}%
  \setlength{\topsep}{0.35em}%
  \setlength{\partopsep}{0pt}%
  \setlength{\itemsep}{0.48em}%
  \setlength{\parsep}{0.18em}%
}{\end{enumerate}}
\newenvironment{examsubparts}{%
  \begin{enumerate}%
  \setlength{\topsep}{0.18em}%
  \setlength{\partopsep}{0pt}%
  \setlength{\itemsep}{0.16em}%
  \setlength{\parsep}{0.1em}%
}{\end{enumerate}}
\newenvironment{examinstructions}{%
  \par\begingroup\itshape
}{%
  \par\endgroup\vspace{0.18em}
}
\makeatletter
\renewcommand\section{\@startsection{section}{1}{0pt}%
  {-1.25ex plus -.25ex minus -.1ex}%
  {0.55ex plus .1ex}%
  {\normalfont\large\bfseries}}
\renewcommand{\@maketitle}{%
  \newpage\null\vskip -1.2em%
  {\footnotesize\bfseries
    \href{https://www.ufl.edu/}{University of Florida}, \href{https://math.ufl.edu/}{Department of Mathematics}\par}%
  \vskip 0.18em%
  \tagstructbegin{tag=Title}%
    {\LARGE\bfseries\href{https://gma.math.ufl.edu/exams/topology-first-year/}{Topology First-Year Exam}, May 2025, Part 1\par}%
  \tagstructend%
  \vskip 0.35em%
  \tagmcbegin{artifact}%
    \hrule height 1pt%
  \tagmcend%
  \vskip 0.55em%
}
\makeatother
\title{Topology First-Year Exam, May 2025, Part 1}
\author{}
\date{}
\begin{document}
\maketitle
\thispagestyle{empty}
\begin{examinstructions}
Answer the following problems and show all your work. Support all statements to the best of your ability.
\end{examinstructions}
\begin{examproblems}{0}
\item
Show that for every set~\(X\), there is no function \(f:X\to 2^X\) such that~\(f\) is onto. Note that \(2^X\) is the power set on~\(X\) and is also denoted by \(\mathcal{P}(X)\).
\item
This problem has two parts.
\begin{examsubparts}
\item[\textnormal{(a)}]
Define what it means for~\(\tau\) to be a topology on a set~\(X\).
\item[\textnormal{(b)}]
If \(\{\tau_\alpha\}_{\alpha\in J}\) is a family of topologies on~\(X\), show that \(\bigcap_{\alpha\in J}\tau_\alpha\) is a topology on~\(X\). (Remark: \(\bigcap_{\alpha\in J}\tau_\alpha=\{U\mid U\in\tau_\alpha\ \forall\alpha\in J\}\).)
\end{examsubparts}
\item
Let \(p:X\to Y\).
\begin{examsubparts}
\item[\textnormal{(a)}]
Define what it means for~\(p\) to be a quotient map.
\item[\textnormal{(b)}]
Assume that~\(p\) is continuous. Show that if there is a continuous map \(f:Y\to X\) such that \(p\circ f\) equals the identity map of~\(Y\), then~\(p\) is a quotient map.
\end{examsubparts}
\item
Let \(f,g:X\to Y\) be continuous maps to a Hausdorff space. Prove that the set \(\{x\mid f(x)=g(x)\}\) is a closed subset of~\(X\).
\item
Prove that if~\(X\) is a compact Hausdorff space, then~\(X\) is a regular space.
\end{examproblems}
\begin{examinstructions}
Answer the following with complete definitions or statements or short proofs.
\end{examinstructions}
\begin{examproblems}{5}
\item
Prove that \(\{(a,b]\mid a<b\}\) is a basis for a topology on~\(\mathbb{R}\).
\item
Let \((X,\tau)\) be a topological space. Let \(A\subset X\). Give the definition of the subspace topology on~\(A\).
\item
Suppose that \(f:X\to Y\) is continuous and that \(A\subset X\) is connected. Show that \(f(A)\subset Y\) is connected.
\item
Does a connected space need to be path connected?
\item
State the Intermediate Value Theorem.
\item
Describe the one-point compactification of \((0,1)\cup(2,3)\).
\item
State the Urysohn Lemma.
\item
Show that a connected normal space~\(X\) having more than one point is uncountable.
\item
What does it mean for a topological space~\(X\) to be a Baire space?
\item
Show that if \(A\subset\mathbb{R}\) is connected then~\(A\) is an interval.
\end{examproblems}
\end{document}
